\begin{abstract}
Aerial person detection supports search and rescue after floods,
wildfires, storms, and at night, yet the benchmarks used to validate it
were collected in calm, clear, well-lit conditions. The operating
conditions that trigger a deployment are precisely the ones that have
never been measured. We present AegisBench, an evaluation benchmark that
places aerial search-and-rescue person detection under nine physically
modelled visual corruptions spanning four disaster families, each at
three severities calibrated against a declared, machine-verified image
statistic. Corruptions are derived from documented atmospheric optics,
including the Koschmieder scattering model for smoke and dust and a
linear-domain photometric model for low light, rather than from generic
image filters. We evaluate three architecturally distinct detectors
across 168 conditions on two datasets spanning the high- and
low-altitude capture regimes, scoring every model with one shared
evaluator at operating points frozen on clean validation data, as a
fielded system would be. Degradation is highly uneven and does not
transfer across capture regimes: heavy rain costs roughly 3.6 points of
recall at the highest severity on one dataset while ranking among the
most damaging corruptions on the other. Low light is universal and
catastrophic, driving recall to exactly zero for all three architectures
on both datasets, with bootstrap confidence intervals degenerate at
$[0.000, 0.000]$; a threshold-independent mAP analysis confirms the
detections are absent rather than below threshold, a three-seed repeat
reproduces the zero exactly, and a check against real night drone
imagery reproduces the collapse. Surviving detections remain accurately
localised, so the failure mode is a detector going blind rather than
mislocalising, which is the mode least visible to a human reviewing the
feed. Finally, we show the collapse is addressable but not free:
fine-tuning on low-light severities 1 and 2 alone lifts severity-3
recall from exactly zero to 0.059--0.188 at a held-out severity, and the
cost to the other eight corruptions is architecture-dependent rather
than uniform, ranging from a net gain to a mean loss of 4.3 points. We
release the corruption engine, the evaluation pipeline, and per-condition
results with intervals for all 168 conditions.
\end{abstract}
